
\documentclass{article}
\usepackage{colortbl}
\usepackage{makecell}
\usepackage{multirow}
\usepackage{supertabular}

\begin{document}

\newcounter{utterance}

\centering \large Interaction Transcript for game `wordle', experiment `high\_frequency\_words\_no\_clue\_no\_critic', episode 3 with reflexion{-}Qwen3.5{-}9B.
\vspace{24pt}

{ \footnotesize  \setcounter{utterance}{1}
\setlength{\tabcolsep}{0pt}
\begin{supertabular}{c@{$\;$}|p{.15\linewidth}@{}p{.15\linewidth}p{.15\linewidth}p{.15\linewidth}p{.15\linewidth}p{.15\linewidth}}
    \# & \multicolumn{2}{c}{Player} && \multicolumn{2}{c}{Game Master} \\
    \hline

    \theutterance \stepcounter{utterance}  
    & & & \multicolumn{4}{p{0.6\linewidth}}{
        \cellcolor[rgb]{0.9,0.9,0.9}{
            \makecell[{{p{\linewidth}}}]{
                \texttt{\tiny{[P1$\langle$GM]}}
                \texttt{You are a language wizard who likes to guess words by using the given rules.} \\
\\ 
\texttt{Welcome to Wordle! You have six attempts to guess the target word, a valid English word of five lowercase letters (a{-}z). Please use the tags "explanation:" and "guess:" to provide a concise explanation for each guess.} \\
\\ 
\texttt{For instance, if your guess is "apple", your response should be} \\
\texttt{explanation: this is a common five{-}letter English word, and I am starting my guess with this word.} \\
\texttt{guess: apple} \\
\\ 
\texttt{After each guess, your answer will be validated, and you will receive feedback indicating which letters are correct (green), which letters are correct but in the wrong position (yellow), and which letters are incorrect (red). This feedback can be useful in determining which letters to include or exclude in your next guess.} \\
\\ 
\texttt{For example, the feedback for "apple" might be:} \\
\texttt{guess\_feedback: a<yellow> p<yellow> p<green> l<yellow> e<red>} \\
\\ 
\texttt{The explanation should contain details about how the guess\_feedback is used to arrive at a new guess.} \\
\\ 
\texttt{Do not generate any other text.} \\
\\ 
\texttt{Let's begin with your first guess.} \\
            }
        }
    }
     \\ \\

    \theutterance \stepcounter{utterance}  
    & \multicolumn{4}{p{0.6\linewidth}}{
        \cellcolor[rgb]{0.561,0.737,0.561}{
            \makecell[{{p{\linewidth}}}]{
                \texttt{\tiny{[P1$\rangle$GM]}}
                \texttt{explanation: I am starting with "crane" because it contains common vowels (a, e) and frequent consonants (r, n) that appear in many five{-}letter words, allowing me to quickly eliminate or confirm several letters.} \\
\texttt{guess: crane} \\
            }
        }
    }
    & & \\ \\

    \theutterance \stepcounter{utterance}  
    & & & \multicolumn{2}{p{0.3\linewidth}}{
        \cellcolor[rgb]{0.9,0.9,0.9}{
            \makecell[{{p{\linewidth}}}]{
                \texttt{\tiny{[GM$|$GM]}}
                \texttt{attempts = 1} \\
\texttt{target\_word = black} \\
\texttt{guess = crane} \\
\texttt{guess\_feedback = c<yellow> r<red> a<green> n<red> e<red>} \\
            }
        }
    }
    & & \\ \\

    \theutterance \stepcounter{utterance}  
    & & & \multicolumn{4}{p{0.6\linewidth}}{
        \cellcolor[rgb]{0.9,0.9,0.9}{
            \makecell[{{p{\linewidth}}}]{
                \texttt{\tiny{[P1$\langle$GM]}}
                \texttt{guess\_feedback: c<yellow> r<red> a<green> n<red> e<red>} \\
\\ 
\texttt{Provide your response only in this format.} \\
\texttt{explanation: details} \\
\texttt{guess: word} \\
            }
        }
    }
     \\ \\

    \theutterance \stepcounter{utterance}  
    & \multicolumn{4}{p{0.6\linewidth}}{
        \cellcolor[rgb]{0.561,0.737,0.561}{
            \makecell[{{p{\linewidth}}}]{
                \texttt{\tiny{[P1$\rangle$GM]}}
                \texttt{explanation: The feedback 'c<yellow> r<red> a<green> n<red> e<red>' indicates that 'a' is the correct letter in the third position, 'c' is present but in the wrong position (likely first or second), and 'r', 'n', 'e' are incorrect. Since 'c' is yellow, it must be in the first or second slot, but not third (occupied by 'a'). Given 'a' is fixed at index 3, and 'c' is not at index 0 (since it was yellow in a guess where it was at 0, wait—if 'c' was at 0 and got yellow, it means it's in the word but not at 0. So 'c' is at index 1 or 2. But index 3 is 'a'. So 'c' is at 1 or 2. Let's try placing 'c' at index 1 and 'a' at index 3. We need a word with pattern ?c?a?. Common words: "acorn" (a at 0, no), "cacao" (c at 0, no), "cocoa" (c at 0, no). Wait, if 'c' is yellow at position 0, it means 'c' is in the word but not at position 0. So 'c' is at 1, 2, or 4. But 'a' is at 3. So 'c' is at 1, 2, or 4. Let's try "cairn" (c at 0, no). How about "acrid"? a at 0, no. Let's try a word where 'c' is at 1 or 2 or 4, 'a' at 3. Example: "cater" (c at 0, no). "cane" is 4 letters. "crane" was the guess. Feedback: c(yellow), r(red), a(green), n(red), e(red). So 'c' is in the word, not at 0. 'a' is at 3. 'r','n','e' are out. So the word has 'c' and 'a' (at 3). Possible words: "cacao" (c at 0, no), "cocoa" (c at 0, no). What about "acorn"? a at 0, no. "caulk"? c at 0, no. "caval"? c at 0, no. Maybe 'c' is at 1? "acids"? a at 0, no. "acids" has a at 0. We need a at 3. Pattern: ? ? a ? ?. And 'c' is in the word but not at} \\
            }
        }
    }
    & & \\ \\

    \theutterance \stepcounter{utterance}  
    & & & \multicolumn{2}{p{0.3\linewidth}}{
        \cellcolor[rgb]{0.9,0.9,0.9}{
            \makecell[{{p{\linewidth}}}]{
                \texttt{\tiny{[GM$|$GM]}}
                \texttt{Error: The response should contain the 'guess:' keyword exactly once.} \\
            }
        }
    }
    & & \\ \\

    \theutterance \stepcounter{utterance}  
    & & & \multicolumn{2}{p{0.3\linewidth}}{
        \cellcolor[rgb]{0.9,0.9,0.9}{
            \makecell[{{p{\linewidth}}}]{
                \texttt{\tiny{[GM$|$GM]}}
                \texttt{game\_result = ABORT} \\
            }
        }
    }
    & & \\ \\

    \theutterance \stepcounter{utterance}  
    & & \multicolumn{4}{p{0.6\linewidth}}{
        \cellcolor[rgb]{0.95,0.95,0.95}{
            \makecell[{{p{\linewidth}}}]{
                \texttt{\tiny{[P1$\leftrightarrow$P]}}
                \texttt{Post{-}episode Reflexion (not sent to GM):} \\
\texttt{New plan: I failed because I generated a long, rambling explanation that did not follow the required format and did not provide a valid guess. The rules strictly require the response to contain only the tags "explanation:" followed by a concise summary of the logic, and "guess:" followed by a valid five{-}letter word. I must stop overthinking the letter placement in the text and instead immediately construct a valid candidate word based on the known constraints (pattern: ? ? a ? ?, with 'c' in positions 1, 2, or 4, and excluding r, n, e). For the next turn, I will select a specific word like "caulk" (if valid) or "cacao" (checking constraints) or better yet, a word fitting ?c?a? or ?ca? or ?c??. Actually, looking at the feedback: c(yellow), a(green at 3). Pattern: \_ \_ a \_ \_. 'c' is not at 0. 'c' is at 1, 2, or 4. Let's try "cacao" (c at 0, invalid). "acorn" (a at 0, invalid). How about "cairn"? c at 0. "cater"? c at 0. "cane"? 4 letters. "cocoa"? c at 0. "caval"? c at 0. "cavalry"? too long. Maybe 'c' is at 1? Pattern: \_ c a \_ \_. Word: "scare"? s{-}c{-}a{-}r{-}e (r is red, invalid). "trace"? t{-}r{-}a{-}c{-}e (r, e invalid). "place"? p{-}l{-}a{-}c{-}e (e invalid). "space"? s{-}p{-}a{-}c{-}e (e invalid). "grace"? g{-}r{-}a{-}c{-}e (r, e invalid). "brace"? b{-}r{-}a{-}c{-}e (r, e invalid). "crane" was the guess. Maybe 'c' is at 2? Pattern: \_ \_ a c \_. "baker"? b{-}a{-}k{-}e{-}r (a at 1, invalid). "maker"? m{-}a{-}k{-}e{-}r (a at 1). "laker"? l{-}a{-}k{-}e{-}r. "saker"? s{-}a{-}k{-}e{-}r. "taker"? t{-}a{-}k{-}e{-}r. "waker"? w{-}a{-}k{-}e{-}r. "baker" has a at 1. We need a at 3. Pattern: \_ \_ a c \_. Word: "baker" no. "laced"? l{-}a{-}c{-}e{-}d (a at 1). "paced"? p{-}a{-}c{-}e{-}d. "raced"? r{-}a{-}c{-}e{-}d. "saced"? "taced"? "vaced"? "waced"? "yaced"? "zaced"? None common. Maybe 'c' is at 4? Pattern} \\
            }
        }
    }
    & \\ \\

\end{supertabular}
}

\end{document}
